\begin{frame}{실습: $\lambda$를 바꿔가며 최적 지점 찾기}
  \begin{itemize}
    \item $\lambda$도 하이퍼파라미터 $\Rightarrow$ $n$과 똑같이
      \textbf{데이터로 고름}.
    \item 같은 random walk, 10 에피소드, $\alpha=0.2$, $\gamma=1$,
      $\lambda$ 0$\sim$1, \textbf{20회 반복 평균} RMS:
  \end{itemize}
  \vskip0.4em
  \small
  \begin{tabular}{@{}ll@{}}
    \toprule
    $\lambda$ & 평균 RMS 오차 \\
    \midrule
    0.0 (TD(0)와 동일) & 0.3172 \\
    0.2 & 0.2790 \\
    0.4 & 0.2329 \\
    0.6 & 0.1774 \\
    \textbf{0.8} & \textbf{0.1377} $\leftarrow$ 최적 \\
    1.0 (사실상 MC) & 0.2674 \\
    \bottomrule
  \end{tabular}
  \vskip0.8em
  \begin{itemize}
    \item 다시 \textbf{U자}. 좌측 $\lambda=0$은 정확히 TD(0)($n=1$,
      $\alpha=0.2$ 결과와 같은 0.3172)라 편향 씻어내지 못함.
      우측 $\lambda=1$은 MC에 가까워 분산 큼.
    \item \textbf{``$\lambda=0.8$이 만능 정답''이 아니다} ---
      환경·에피소드 수(예: 1000)가 바뀌면 U자 바닥은 \textbf{완전히
      다른 위치}로 이동. 에피소드가 충분히 많으면 편향이 씻겨
      $\lambda$를 1에 가깝게(MC 쪽) 가져가도 좋은 경우 많음.
    \item \kq{``데이터로 고르는 것'' 자체가 방법 --- 검증 데이터·
      반복 실험으로 매번 재확인하는 습관이 실전의 본체.}
  \end{itemize}
\end{frame}
